% Einleitung: 1/2-3/4 Seite, mit konkreten Beiträgen (Folie 16).
\section{Einleitung}
\label{sec:intro}

Auf großen Datensätzen vortrainierte Vision Transformer (ViTs)~\cite{dosovitskiy2021vit} sind starke
Ausgangspunkte für neue Bildklassifikationsaufgaben. Bei kleinen Zieldatensätzen stellt sich die Frage,
\emph{wie} man sie anpasst: Nur einen neuen Klassifikator auf eingefrorenen Merkmalen zu trainieren
(Linear Probing) ist günstig, trennt die Zielklassen aber womöglich nicht gut genug. Alle Gewichte zu
aktualisieren (volles Fine-Tuning) ist teuer und kann die vortrainierten Merkmale
verzerren~\cite{kumar2022finetuning}. Low-Rank Adaptation (LoRA)~\cite{hu2022lora} bietet einen
Mittelweg: Die vortrainierten Gewichte bleiben eingefroren, gelernt werden nur kleine Änderungen niedrigen
Rangs.

Wir untersuchen diesen Zielkonflikt an einer feinkörnigen Aufgabe. Der Empfehlung folgend, nur eine
Tierart zu verwenden, klassifizieren wir die 25 Hunderassen von Oxford-IIIT Pet~\cite{parkhi2012pets};
Rassen derselben Art sehen sich ähnlich und unterscheiden sich in Details von Kopfform, Fell und
Körperproportionen. Die Situation ist typisch für die Praxis: wenige gelabelte Bilder (etwa 90 pro
Rasse für das Training), viele ähnliche Klassen und begrenzte Rechenleistung.

\paragraph{Forschungsfrage.}
Wie effizient lässt sich ein vortrainierter ViT an eine neue Aufgabe anpassen -- welche Accuracy
erreichen Linear Probing und LoRA im Vergleich zu vollem Fine-Tuning und zu CNNs, und zu welchen Kosten?
Wir prüfen vier Hypothesen:
\textbf{A1} LoRA mit $<1\,\%$ trainierbaren Parametern liegt weniger als 2 Prozentpunkte unter vollem
Fine-Tuning;
\textbf{A2} Linear Probing ist bei feinkörnigen Rassen deutlich schlechter als LoRA;
\textbf{A3} ein höherer LoRA-Rang als $r=8$ bringt wenig, mehr angepasste Layer bringen mehr;
\textbf{A4} der beste ViT und die CNNs erreichen eine ähnliche Accuracy.

\paragraph{Beiträge.}
\begin{itemize}
  \item Ein Vergleich von Linear Probing, LoRA und vollem Fine-Tuning des ViT-B/16 mit ResNet-50 und
        ConvNeXt-T unter identischem Protokoll (drei Seeds, Modellauswahl nur auf Validation).
  \item Eine LoRA-Ablation über Rang und angepasste Layer.
  \item Eine Kosten-Nutzen-Analyse: Accuracy gegenüber trainierbaren Parametern und Trainingszeit sowie
        der Einfluss der Datenaugmentierung.
  \item Eine Analyse, \emph{warum} alle Methoden über 90\,\% erreichen (Training ohne Pretraining,
        Überschneidung mit den ImageNet-Klassen, Suche nach Fast-Duplikaten), und Erklärungskarten,
        worauf die Modelle schauen (\cref{sec:why,sec:explain}).
  \item Eine Datenqualitätsanalyse der 25 Hunderassen sowie dokumentierter, reproduzierbarer Code
        (\url{https://github.com/hsbi-dlcv-eskeif-dawalibi/hunderassen-vit}).
\end{itemize}
